\documentclass[aspectratio=169]{beamer}
\usetheme{Madrid}
\usecolortheme{default}

% Input the automated results summary macros (go up one directory to report/)
%% Automated results summary generated by generate_latex_results.py
%% Do not edit manually!

\newcommand{\pretrainedfrozenValEpochsMean}{8.0}
\newcommand{\pretrainedfrozenValEpochsStd}{0.8}
\newcommand{\pretrainedfrozenValAccMean}{93.82\%}
\newcommand{\pretrainedfrozenValAccStd}{0.32\%}
\newcommand{\pretrainedfrozenTestAccMean}{93.27\%}
\newcommand{\pretrainedfrozenTestAccStd}{0.23\%}
\newcommand{\pretrainedfrozenClosestSeed}{100}
\newcommand{\pretrainedfrozenRepPrecBuildings}{93.2\%}
\newcommand{\pretrainedfrozenRepRecBuildings}{93.4\%}
\newcommand{\pretrainedfrozenRepFBuildings}{93.3\%}
\newcommand{\pretrainedfrozenRepPrecForest}{98.7\%}
\newcommand{\pretrainedfrozenRepRecForest}{99.4\%}
\newcommand{\pretrainedfrozenRepFForest}{99.1\%}
\newcommand{\pretrainedfrozenRepPrecGlacier}{93.9\%}
\newcommand{\pretrainedfrozenRepRecGlacier}{82.8\%}
\newcommand{\pretrainedfrozenRepFGlacier}{88.0\%}
\newcommand{\pretrainedfrozenRepPrecMountain}{86.3\%}
\newcommand{\pretrainedfrozenRepRecMountain}{92.2\%}
\newcommand{\pretrainedfrozenRepFMountain}{89.1\%}
\newcommand{\pretrainedfrozenRepPrecSea}{95.1\%}
\newcommand{\pretrainedfrozenRepRecSea}{98.6\%}
\newcommand{\pretrainedfrozenRepFSea}{96.8\%}
\newcommand{\pretrainedfrozenRepPrecStreet}{93.3\%}
\newcommand{\pretrainedfrozenRepRecStreet}{94.4\%}
\newcommand{\pretrainedfrozenRepFStreet}{93.8\%}
\newcommand{\pretrainedfrozenRepAcc}{93.23\%}
\newcommand{\pretrainedfrozenCM}[2]{%
  \ifnum#1=0\relax\ifnum#2=0\relax{408}\fi\fi%
  \ifnum#1=0\relax\ifnum#2=1\relax{0}\fi\fi%
  \ifnum#1=0\relax\ifnum#2=2\relax{0}\fi\fi%
  \ifnum#1=0\relax\ifnum#2=3\relax{0}\fi\fi%
  \ifnum#1=0\relax\ifnum#2=4\relax{1}\fi\fi%
  \ifnum#1=0\relax\ifnum#2=5\relax{28}\fi\fi%
  \ifnum#1=1\relax\ifnum#2=0\relax{0}\fi\fi%
  \ifnum#1=1\relax\ifnum#2=1\relax{471}\fi\fi%
  \ifnum#1=1\relax\ifnum#2=2\relax{0}\fi\fi%
  \ifnum#1=1\relax\ifnum#2=3\relax{1}\fi\fi%
  \ifnum#1=1\relax\ifnum#2=4\relax{0}\fi\fi%
  \ifnum#1=1\relax\ifnum#2=5\relax{2}\fi\fi%
  \ifnum#1=2\relax\ifnum#2=0\relax{1}\fi\fi%
  \ifnum#1=2\relax\ifnum#2=1\relax{4}\fi\fi%
  \ifnum#1=2\relax\ifnum#2=2\relax{458}\fi\fi%
  \ifnum#1=2\relax\ifnum#2=3\relax{71}\fi\fi%
  \ifnum#1=2\relax\ifnum#2=4\relax{16}\fi\fi%
  \ifnum#1=2\relax\ifnum#2=5\relax{3}\fi\fi%
  \ifnum#1=3\relax\ifnum#2=0\relax{1}\fi\fi%
  \ifnum#1=3\relax\ifnum#2=1\relax{2}\fi\fi%
  \ifnum#1=3\relax\ifnum#2=2\relax{28}\fi\fi%
  \ifnum#1=3\relax\ifnum#2=3\relax{484}\fi\fi%
  \ifnum#1=3\relax\ifnum#2=4\relax{9}\fi\fi%
  \ifnum#1=3\relax\ifnum#2=5\relax{1}\fi\fi%
  \ifnum#1=4\relax\ifnum#2=0\relax{2}\fi\fi%
  \ifnum#1=4\relax\ifnum#2=1\relax{0}\fi\fi%
  \ifnum#1=4\relax\ifnum#2=2\relax{2}\fi\fi%
  \ifnum#1=4\relax\ifnum#2=3\relax{3}\fi\fi%
  \ifnum#1=4\relax\ifnum#2=4\relax{503}\fi\fi%
  \ifnum#1=4\relax\ifnum#2=5\relax{0}\fi\fi%
  \ifnum#1=5\relax\ifnum#2=0\relax{26}\fi\fi%
  \ifnum#1=5\relax\ifnum#2=1\relax{0}\fi\fi%
  \ifnum#1=5\relax\ifnum#2=2\relax{0}\fi\fi%
  \ifnum#1=5\relax\ifnum#2=3\relax{2}\fi\fi%
  \ifnum#1=5\relax\ifnum#2=4\relax{0}\fi\fi%
  \ifnum#1=5\relax\ifnum#2=5\relax{473}\fi\fi%
}
\newcommand{\randomfrozenValEpochsMean}{8.7}
\newcommand{\randomfrozenValEpochsStd}{2.1}
\newcommand{\randomfrozenValAccMean}{67.83\%}
\newcommand{\randomfrozenValAccStd}{4.62\%}
\newcommand{\randomfrozenTestAccMean}{66.54\%}
\newcommand{\randomfrozenTestAccStd}{4.04\%}
\newcommand{\randomfrozenClosestSeed}{100}
\newcommand{\randomfrozenRepPrecBuildings}{53.0\%}
\newcommand{\randomfrozenRepRecBuildings}{64.5\%}
\newcommand{\randomfrozenRepFBuildings}{58.2\%}
\newcommand{\randomfrozenRepPrecForest}{62.5\%}
\newcommand{\randomfrozenRepRecForest}{97.9\%}
\newcommand{\randomfrozenRepFForest}{76.3\%}
\newcommand{\randomfrozenRepPrecGlacier}{78.7\%}
\newcommand{\randomfrozenRepRecGlacier}{53.5\%}
\newcommand{\randomfrozenRepFGlacier}{63.7\%}
\newcommand{\randomfrozenRepPrecMountain}{67.1\%}
\newcommand{\randomfrozenRepRecMountain}{69.5\%}
\newcommand{\randomfrozenRepFMountain}{68.3\%}
\newcommand{\randomfrozenRepPrecSea}{76.4\%}
\newcommand{\randomfrozenRepRecSea}{49.4\%}
\newcommand{\randomfrozenRepFSea}{60.0\%}
\newcommand{\randomfrozenRepPrecStreet}{68.1\%}
\newcommand{\randomfrozenRepRecStreet}{64.7\%}
\newcommand{\randomfrozenRepFStreet}{66.3\%}
\newcommand{\randomfrozenRepAcc}{66.10\%}
\newcommand{\randomfrozenCM}[2]{%
  \ifnum#1=0\relax\ifnum#2=0\relax{282}\fi\fi%
  \ifnum#1=0\relax\ifnum#2=1\relax{64}\fi\fi%
  \ifnum#1=0\relax\ifnum#2=2\relax{0}\fi\fi%
  \ifnum#1=0\relax\ifnum#2=3\relax{11}\fi\fi%
  \ifnum#1=0\relax\ifnum#2=4\relax{7}\fi\fi%
  \ifnum#1=0\relax\ifnum#2=5\relax{73}\fi\fi%
  \ifnum#1=1\relax\ifnum#2=0\relax{4}\fi\fi%
  \ifnum#1=1\relax\ifnum#2=1\relax{464}\fi\fi%
  \ifnum#1=1\relax\ifnum#2=2\relax{0}\fi\fi%
  \ifnum#1=1\relax\ifnum#2=3\relax{1}\fi\fi%
  \ifnum#1=1\relax\ifnum#2=4\relax{0}\fi\fi%
  \ifnum#1=1\relax\ifnum#2=5\relax{5}\fi\fi%
  \ifnum#1=2\relax\ifnum#2=0\relax{64}\fi\fi%
  \ifnum#1=2\relax\ifnum#2=1\relax{29}\fi\fi%
  \ifnum#1=2\relax\ifnum#2=2\relax{296}\fi\fi%
  \ifnum#1=2\relax\ifnum#2=3\relax{86}\fi\fi%
  \ifnum#1=2\relax\ifnum#2=4\relax{33}\fi\fi%
  \ifnum#1=2\relax\ifnum#2=5\relax{45}\fi\fi%
  \ifnum#1=3\relax\ifnum#2=0\relax{41}\fi\fi%
  \ifnum#1=3\relax\ifnum#2=1\relax{41}\fi\fi%
  \ifnum#1=3\relax\ifnum#2=2\relax{36}\fi\fi%
  \ifnum#1=3\relax\ifnum#2=3\relax{365}\fi\fi%
  \ifnum#1=3\relax\ifnum#2=4\relax{36}\fi\fi%
  \ifnum#1=3\relax\ifnum#2=5\relax{6}\fi\fi%
  \ifnum#1=4\relax\ifnum#2=0\relax{89}\fi\fi%
  \ifnum#1=4\relax\ifnum#2=1\relax{26}\fi\fi%
  \ifnum#1=4\relax\ifnum#2=2\relax{43}\fi\fi%
  \ifnum#1=4\relax\ifnum#2=3\relax{77}\fi\fi%
  \ifnum#1=4\relax\ifnum#2=4\relax{252}\fi\fi%
  \ifnum#1=4\relax\ifnum#2=5\relax{23}\fi\fi%
  \ifnum#1=5\relax\ifnum#2=0\relax{52}\fi\fi%
  \ifnum#1=5\relax\ifnum#2=1\relax{118}\fi\fi%
  \ifnum#1=5\relax\ifnum#2=2\relax{1}\fi\fi%
  \ifnum#1=5\relax\ifnum#2=3\relax{4}\fi\fi%
  \ifnum#1=5\relax\ifnum#2=4\relax{2}\fi\fi%
  \ifnum#1=5\relax\ifnum#2=5\relax{324}\fi\fi%
}
\newcommand{\randomfullValEpochsMean}{15.0}
\newcommand{\randomfullValEpochsStd}{2.2}
\newcommand{\randomfullValAccMean}{86.31\%}
\newcommand{\randomfullValAccStd}{1.77\%}
\newcommand{\randomfullTestAccMean}{84.69\%}
\newcommand{\randomfullTestAccStd}{1.87\%}
\newcommand{\randomfullClosestSeed}{100}
\newcommand{\randomfullRepPrecBuildings}{82.5\%}
\newcommand{\randomfullRepRecBuildings}{82.8\%}
\newcommand{\randomfullRepFBuildings}{82.6\%}
\newcommand{\randomfullRepPrecForest}{96.4\%}
\newcommand{\randomfullRepRecForest}{95.6\%}
\newcommand{\randomfullRepFForest}{96.0\%}
\newcommand{\randomfullRepPrecGlacier}{89.2\%}
\newcommand{\randomfullRepRecGlacier}{70.2\%}
\newcommand{\randomfullRepFGlacier}{78.5\%}
\newcommand{\randomfullRepPrecMountain}{76.4\%}
\newcommand{\randomfullRepRecMountain}{85.0\%}
\newcommand{\randomfullRepFMountain}{80.4\%}
\newcommand{\randomfullRepPrecSea}{84.7\%}
\newcommand{\randomfullRepRecSea}{88.2\%}
\newcommand{\randomfullRepFSea}{86.5\%}
\newcommand{\randomfullRepPrecStreet}{82.8\%}
\newcommand{\randomfullRepRecStreet}{89.4\%}
\newcommand{\randomfullRepFStreet}{86.0\%}
\newcommand{\randomfullRepAcc}{84.90\%}
\newcommand{\randomfullCM}[2]{%
  \ifnum#1=0\relax\ifnum#2=0\relax{362}\fi\fi%
  \ifnum#1=0\relax\ifnum#2=1\relax{3}\fi\fi%
  \ifnum#1=0\relax\ifnum#2=2\relax{2}\fi\fi%
  \ifnum#1=0\relax\ifnum#2=3\relax{1}\fi\fi%
  \ifnum#1=0\relax\ifnum#2=4\relax{6}\fi\fi%
  \ifnum#1=0\relax\ifnum#2=5\relax{63}\fi\fi%
  \ifnum#1=1\relax\ifnum#2=0\relax{2}\fi\fi%
  \ifnum#1=1\relax\ifnum#2=1\relax{453}\fi\fi%
  \ifnum#1=1\relax\ifnum#2=2\relax{0}\fi\fi%
  \ifnum#1=1\relax\ifnum#2=3\relax{4}\fi\fi%
  \ifnum#1=1\relax\ifnum#2=4\relax{2}\fi\fi%
  \ifnum#1=1\relax\ifnum#2=5\relax{13}\fi\fi%
  \ifnum#1=2\relax\ifnum#2=0\relax{6}\fi\fi%
  \ifnum#1=2\relax\ifnum#2=1\relax{5}\fi\fi%
  \ifnum#1=2\relax\ifnum#2=2\relax{388}\fi\fi%
  \ifnum#1=2\relax\ifnum#2=3\relax{111}\fi\fi%
  \ifnum#1=2\relax\ifnum#2=4\relax{36}\fi\fi%
  \ifnum#1=2\relax\ifnum#2=5\relax{7}\fi\fi%
  \ifnum#1=3\relax\ifnum#2=0\relax{12}\fi\fi%
  \ifnum#1=3\relax\ifnum#2=1\relax{3}\fi\fi%
  \ifnum#1=3\relax\ifnum#2=2\relax{30}\fi\fi%
  \ifnum#1=3\relax\ifnum#2=3\relax{446}\fi\fi%
  \ifnum#1=3\relax\ifnum#2=4\relax{32}\fi\fi%
  \ifnum#1=3\relax\ifnum#2=5\relax{2}\fi\fi%
  \ifnum#1=4\relax\ifnum#2=0\relax{15}\fi\fi%
  \ifnum#1=4\relax\ifnum#2=1\relax{3}\fi\fi%
  \ifnum#1=4\relax\ifnum#2=2\relax{14}\fi\fi%
  \ifnum#1=4\relax\ifnum#2=3\relax{20}\fi\fi%
  \ifnum#1=4\relax\ifnum#2=4\relax{450}\fi\fi%
  \ifnum#1=4\relax\ifnum#2=5\relax{8}\fi\fi%
  \ifnum#1=5\relax\ifnum#2=0\relax{42}\fi\fi%
  \ifnum#1=5\relax\ifnum#2=1\relax{3}\fi\fi%
  \ifnum#1=5\relax\ifnum#2=2\relax{1}\fi\fi%
  \ifnum#1=5\relax\ifnum#2=3\relax{2}\fi\fi%
  \ifnum#1=5\relax\ifnum#2=4\relax{5}\fi\fi%
  \ifnum#1=5\relax\ifnum#2=5\relax{448}\fi\fi%
}


\title[Transfer Learning vs. Scratch]{Transfer Learning versus Training from Scratch for Scene Image Classification}
\subtitle{CS 5330 Final Project}
\author{Shriman Raghav Srinivasan}
\institute[Northeastern]{Khoury College of Computer Sciences \\ Northeastern University}
\date{July 2026}

\begin{document}

% Slide 1: Title
\begin{frame}
  \titlepage
\end{frame}

% Slide 2: Introduction & Motivation
\begin{frame}{Introduction \& Motivation}
  \begin{block}{The Core Question}
    How much of the benefit of Transfer Learning is due to \textbf{pre-trained feature quality}, and how much is due to \textbf{reduced trainable capacity} (fewer parameters to optimize)?
  \end{block}
  \begin{itemize}
    \item A naive comparison between a pre-trained frozen network and a fully trained random network confounds feature quality and parameter capacity.
    \item \textbf{Our Solution:} A capacity-matched three-condition experimental design:
      \begin{enumerate}
        \item \textbf{Pretrained-Frozen (PT-F):} Loaded with ImageNet weights; early layers frozen.
        \item \textbf{Random-Frozen (R-F):} Randomly initialized; early layers frozen (Capacity Control).
        \item \textbf{Random-Full (R-FL):} Randomly initialized; all layers trainable (Scratch Baseline).
      \end{enumerate}
  \end{itemize}
\end{frame}

% Slide 3: Model Architecture \& Partitioning
\begin{frame}{Model Architecture \& Partitioning}
  \begin{itemize}
    \item \textbf{Base Model:} ResNet-18 architecture.
    \item \textbf{Frozen Subset ($\theta_f$):} Input convolution (\texttt{conv1}), first batchnorm (\texttt{bn1}), and the first three residual stages (\texttt{layer1}, \texttt{layer2}, \texttt{layer3}) contain 2,782,784 parameters.
    \item \textbf{Trainable Subset ($\theta_t$):} Final residual stage (\texttt{layer4}) and 6-class output head (\texttt{fc}) contain 8,396,806 parameters.
    \item \textbf{Critical Technical Detail:} Batch Normalization layers in $\theta_f$ are forced into evaluation mode (\texttt{model.eval()}) during training. This prevents running statistics from drifting and ruining the representation.
  \end{itemize}
\end{frame}

% Slide 4: Dataset \& Preprocessing
\begin{frame}{Dataset \& Preprocessing}
  \begin{columns}
    \begin{column}{0.5\textwidth}
      \textbf{Intel Image Classification Dataset}
      \begin{itemize}
        \item Natural scene images across 6 classes: \textit{Buildings, Forest, Glacier, Mountain, Sea, Street}.
        \item Preprocessing:
          \begin{enumerate}
            \item Resized shorter side to 224px.
            \item Center-cropped to 224x224.
            \item ImageNet channel mean \& std dev normalization.
          \end{enumerate}
      \end{itemize}
    \end{column}
    \begin{column}{0.5\textwidth}
      \textbf{Splits Structure}
      \begin{itemize}
        \item \textbf{Training (90\%):} 12,631 images.
        \item \textbf{Validation (10\%):} 1,403 images (reproducible class-stratified split).
        \item \textbf{Test (Held-out):} 3,000 images (exclusively for final reporting).
      \end{itemize}
    \end{column}
  \end{columns}
\end{frame}

% Slide 5: Experimental Configuration
\begin{frame}{Experimental Configuration}
  \begin{block}{Optimizer \& Loss}
    We optimize the trainable parameters $\theta_t$ using Stochastic Gradient Descent (SGD) with Cross-Entropy Loss.
  \end{block}
  \begin{itemize}
    \item \textbf{Hyperparameters:} Learning rate = 0.001, momentum = 0.9, weight decay = 1e-4. Batch size = 32.
    \item \textbf{Early Stopping:} Validation loss is monitored; training stops after 5 consecutive epochs of no improvement.
    \item \textbf{Robustness:} Each of the 3 conditions runs across 3 distinct random seeds (42, 100, 2026) to generate statistics.
  \end{itemize}
\end{frame}

% Slide 6: Results - Accuracy and Convergence
\begin{frame}{Results - Accuracy and Convergence}
  \begin{table}
    \centering
    \caption{Test Accuracy and Convergence Speed Comparison}
    \begin{tabular}{lccc}
      \hline
      \textbf{Condition} & \textbf{Mean Test Acc.} & \textbf{Std. Dev.} & \textbf{Mean Epochs} \\ \hline
      Pretrained-Frozen & \pretrainedfrozenTestAccMean & \pretrainedfrozenTestAccStd & \pretrainedfrozenValEpochsMean \\
      Random-Frozen & \randomfrozenTestAccMean & \randomfrozenTestAccStd & \randomfrozenValEpochsMean \\
      Random-Full & \randomfullTestAccMean & \randomfullTestAccStd & \randomfullValEpochsMean \\ \hline
    \end{tabular}
  \end{table}
  \begin{itemize}
    \item \textbf{Key Takeaway:} PT-F achieves a massive accuracy boost ($>50\%$) over R-F, directly isolating the value of ImageNet pretrained representations.
  \end{itemize}
\end{frame}

% Slide 7: Error Analysis \& Class Confusions
\begin{frame}{Error Analysis \& Class Confusions}
  \begin{itemize}
    \item Representative seeds closest to the mean test accuracy:
      \begin{itemize}
        \item Pretrained-Frozen: Seed \pretrainedfrozenClosestSeed\ (\pretrainedfrozenRepAcc)
        \item Random-Frozen: Seed \randomfrozenClosestSeed\ (\randomfrozenRepAcc)
        \item Random-Full: Seed \randomfullClosestSeed\ (\randomfullRepAcc)
      \end{itemize}
    \item \textbf{Common Confusions Across All Conditions:}
      \begin{itemize}
        \item \textbf{Glacier vs. Mountain:} Visually similar snow layers and rock features.
        \item \textbf{Buildings vs. Street:} Buildings naturally populate street backgrounds.
      \end{itemize}
    \item Confusion is a property of dataset visual similarity, not the initialization or capacity condition.
  \end{itemize}
\end{frame}

% Slide 8: Discussion \& Takeaways
\begin{frame}{Discussion \& Takeaways}
  \begin{itemize}
    \item \textbf{Pretraining is Essential:} For matched capacities, ImageNet pretraining yields a performance boost of $>50\%$. Arbitrary random weights in $\theta_f$ constrain the representation space.
    \item \textbf{Scratch Limits:} Training from scratch with full capacity (R-FL) outperforms the frozen random control (R-F) but remains far inferior to transfer learning (PT-F), despite taking much longer to converge.
    \item \textbf{Final Conclusion:} In compute-constrained and data-limited settings, pre-trained weights function as critical regularizers and convergence accelerators.
  \end{itemize}
\end{frame}

\end{document}
